
\def\PUDcodigo{EME106}

\if\COMPILAR0
    \def\PUDcodacad{04507.11}
\else
    \def\PUDcodacad{04506.10}
\fi

\def\PUDdisciplina{Desenho Auxiliado por Computador}

\def\PUDsemestre{S2}

\def\PUDhoras{80}

%NOVOS ---------------------------------------------------
\def\PUDhorasTeoria{20}
\def\PUDhorasPratica{60}

\if\COMPILAR0
   \def\PUDinterECA{\hyperref[PUD:ACE004]{Inglês Técnico (04507.3)} 
   
   \hyperref[PUD:EME105]{Desenho Técnico (04507.6)} 
   
   \hyperref[PUD:EME123]{Manufatura Auxiliada por Computador (04507.46)}
   }
\else
   \def\PUDinterECA{\hyperref[PUD:ACE004]{Inglês Técnico (04507.3)} 
   
   \hyperref[PUD:EME105]{Desenho Técnico (04506.6)} 
   
   \hyperref[PUD:EME123]{Manufatura Auxiliada por Computador (04506.42)}
   }
\fi

\def\PUDrecursos{Projetor multimídia; Quadro branco e pincel; Aulas em laboratório de informática.}

\def\PUDatuaisECA{---}

%----------------------------------------------------------

\def\PUDcreditos{4}

\def\PUDprerequisito{ \hyperref[PUD:EME105]{EME105} }

\if\COMPILAR0
    \def\PUDpreacad{ \hyperref[PUD:EME105]{Desenho Técnico (04507.6)} }
\else
    \def\PUDpreacad{ \hyperref[PUD:EME105]{Desenho Técnico (04506.6)} }
\fi

\def\PUDobjetivos{Conhecer as características de um sistema CAD e suas aplicações. Estudar os principais comandos do sistema CAD utilizado para representar desenhos em 2D. Construir modelos tridimensionais de maquinas, componentes ou dispositivos mecânicos, montar conjuntos e simular movimentos em ambiente digital.
%Projetar as principais máquinas de levantamento e transporte utilizadas na indústria.
}

\def\PUDementa{Introdução ao ambiente CAD (definições importantes, ambiente de trabalho, unidades de trabalho, símbolos especiais).          Desenho em perspectiva isométrica.          Sistemas de Coordenadas.          Comando de visualização, criação, edição e dimensionamento.         Noções de CAD 3D.}

\def\PUDconteudo{ & Introdução, definição, ambiente, unidade e símbolos utilizados em CAD;  Teclas de atalho, símbolos especiais; 
\\ 
 & Desenho em perspectiva isométrica; 
\\ 
 & Sistemas de Coordenadas; 
\\ 
 & Comandos de visualização, criação, edição e dimensionamento de desenhos em 2D; 
\\ 
 & Noções de CAD 3D. }

\def\PUDmetodo{Aulas expositórias que podem ser teóricas e/ou práticas, onde as práticas no laboratório serão marcadas no decorrer da disciplina.}

\def\PUDavalia{A avaliação será feita com aplicação de provas teóricas e/ou práticas, além da possibilidade de inclusão de trabalhos e seminários no decorrer da disciplina.}

\def\PUDbibbasica{ & \href{www.biblioteca.ifce.edu.br/index.asp?codigo_sophia=24065}{Harrington}, D. J.  Desvendando o AutoCAD 2005.  São Paulo, SP: Pearson Makron Books, 2006.  716p.  ISBN: 8534615446. 
\\ 
 & \href{www.biblioteca.ifce.edu.br/index.asp?codigo_sophia=24059}{Saad}, A. L.  AutoCAD 2004 2D e 3D para Engenharia e Arquitetura.  São Paulo, SP: Pearson Makron Books, 2004.  280p. ISBN: 8534615357. 
\\ 
 & \href{www.biblioteca.ifce.edu.br/index.asp?codigo_sophia=24129}{Lima}, C. C.  Estudo dirigido de AutoCAD 2007.  4º ed.  São Paulo, SP : Érica, 2011.  300p.  ISBN: 9788536501185. }

\def\PUDbibcomplementar{ & \href{www.biblioteca.ifce.edu.br/index.asp?codigo_sophia=23779}{Junghans}, Daniel.  Informática aplicada ao desenho técnico.  Curitiba, PR: Base Editorial, 2010.  224p.  ISBN: 9788579055478. 
\\ 
 &  \href{www.biblioteca.ifce.edu.br/index.asp?codigo_sophia=23872}{Silva}, Arlindo.  Desenho técnico moderno.  4º ed.  Rio de Janeiro, RJ : LTC, 2006. 475p.  ISBN: 9788521615224.
 %& Frey, David.  Autocad 2002: a bíblia do iniciante.  Rio de Janeiro, RJ: Ciência Moderna, 2003.  560p.  ISBN: 8573932414. 
\\ 
 & \href{www.biblioteca.ifce.edu.br/index.asp?codigo_sophia=24198}{Strauhs}, Faimara do Rocio.  Desenho Técnico.  Curitiba, PR: Base Editorial, 2010. 112p.  ISBN: 9788579055393.
\\
 & \href{www.biblioteca.ifce.edu.br/index.asp?codigo_sophia=55358}{Schwartz}, Steve; Davis, Phyllis.  CorelDRAW 11: passo a passo Lite.  E-book.  Pearson.  224p.  ISBN: 9788534615082.
\\
 & \href{www.biblioteca.ifce.edu.br/index.asp?codigo_sophia=24501}{Baldam}, Roquemar de Lima.  AutoCAD 2007: utilizando totalmente.  2º ed.  São Paulo, SP: Érica, 2008.  458p.  ISBN: 9788536501550. }

\def\PUDautor{Samuel Vieira Dias}

\def\PUDdata{2013-04-15}

\def\PUDrevisao{2}

\def\PUDrevisor{José Ciro dos Santos}

\def\PUDdatarevisao{2019-05-15}

\PUDcorpo
